\documentclass[sigconf]{acmart}
\setcopyright{none}
\renewcommand\footnotetextcopyrightpermission[1]{}
\settopmatter{printacmref=false, printfolios=true}
\acmYear{2026}

% acmart loads amsmath, amssymb, hyperref, natbib, xcolor, geometry,
% booktabs, graphicx internally

% Custom commands
\newcommand{\eml}{\mathrm{eml}}
\newcommand{\sml}{\mathrm{sml}}
\newcommand{\rml}{\mathrm{rml}}
\newcommand{\op}{\mathrm{op}}

\title{Why Architecture Choice Matters in Symbolic Regression}

\author{Chakshu Gupta}
\email{cgupta65@gatech.edu}
\affiliation{%
  \institution{Georgia Institute of Technology}
  \city{Khurja}
  \country{India}
}

\begin{document}
\fancyhead{}  % suppress running headers

\begin{abstract}
Symbolic regression discovers mathematical formulas from data.
Some methods fix a tree of operators, assign learnable weights, and
train by gradient descent.  The tree's structure, which determines
what operators and variables appear at each position, is chosen once
and applied to every target.  This paper tests whether that choice
affects which targets are actually recovered.  Three structures are
compared, all sharing the same operator and target language but
differing in how variables enter the tree; one is strictly more
expressive.  Across over 12{,}700 training runs, one structure
recovers a target at 100\% while another scores~0\%, and the ranking
reverses on a different target.  Expressiveness guarantees that a
solution exists in the search space, but not that gradient descent
finds it: the most expressive structure fails on targets that a
restricted alternative solves reliably.  Switching the operator
changes which targets succeed; reversing its gradient profile
collapses recovery entirely.  Balanced (non-chain) tree shapes are
never recovered.  These findings show that the optimization landscape,
not expressiveness alone, determines what gradient-based symbolic
regression recovers.
\end{abstract}

\maketitle


%% ===================================================================
\section{Introduction}
\label{sec:intro}
%% ===================================================================

Symbolic regression (SR) discovers mathematical formulas from
data~\citep{LaCava2021srbench, Dong2025survey, Makke2024survey}.  One
family of methods parameterizes a formula as a tree of operators
whose inputs are controlled by learnable weights, and recovers the
target by training those weights with gradient
descent~\citep{Sahoo2018eql, Dong2024eeql}.  The tree's
structure, which determines what operators and variables are available
at each position, is its \emph{architecture}: a design choice made
once and applied to every target.

Prior work has compared SR methods that differ in architecture,
grammar, and search simultaneously: SRBench~\citep{LaCava2021srbench}
tests 14 methods on 252 problems but cannot isolate architecture
effects.  Gradient-based methods such as
EQL~\citep{Sahoo2018eql} and eEQL~\citep{Dong2024eeql} each use a
single fixed architecture without studying how architecture interacts
with targets.  MetaSymNet~\citep{Li2025metasymnet} meta-learns the
architecture but does not explain why fixed ones fail.
Reinforcement-learning approaches~\citep{Petersen2021dsr,
Landajuela2022udsr, Biggio2021nsrs} avoid fixed-architecture
landscapes entirely.  On the EML side,
Odrzywolek~\citep{Odrzywolek2026eml} demonstrated SR with one
architecture, and Ipek~\citep{Ipek2026hardware, Ipek2026battery}
applied EML to hardware and battery models without interaction
analysis.  No prior study has varied architecture
while holding the target set constant.

This paper isolates the architecture variable by using the EML
(Exp-Minus-Log) operator~\citep{Odrzywolek2026eml}, a single
binary operation from which all elementary functions can be
constructed.  Because every node in an EML tree computes the same
operation, targets differ only in tree shape and variable placement,
making it possible to compare architectures that share the same
target language.  Three architectures are tested across more than
12{,}700 training runs; the results show that a target's
recoverability depends on architecture, not just on whether the
architecture can express it.


%% ===================================================================
\section{Background}
\label{sec:background}
%% ===================================================================

The EML operator~\citep{Odrzywolek2026eml} computes
$\eml(a,b) = e^a - \ln b$.
Its two inputs receive asymmetric gradients: the left input is
exponentially amplified ($\partial\eml/\partial a = e^a$) while
the right is attenuated ($\partial\eml/\partial b = -1/b$).  EML
is a Sheffer operator: with only the constant~$1$, binary trees of
EML nodes can express every elementary function ($+$, $-$, $\times$,
$/$, $\exp$, $\ln$, trigonometric, and
more)~\citep{Odrzywolek2026eml}.  For instance, $e^x = \eml(x,1)$
and $e = \eml(1,1)$.

A depth-$n$ EML tree is a perfect binary tree with $2^n - 1$
operator nodes (depth~3 gives 7~nodes with 8~leaf positions).
At depth~3, five distinct shapes exist.
Four are \emph{chains}, in which one branch at each level carries the
computation while the other is a leaf.  Chains are classified by which
EML input (Left or Right) carries the active branch:

\begin{center}
\begin{tabular}{@{}lll@{}}
\toprule
Shape & Template & Active-branch path \\
\midrule
LR & $\eml(\eml(1,\eml(\cdot,\cdot)),1)$ & Left at root, Right at depth~1 \\
RL & $\eml(1,\eml(\eml(\cdot,\cdot),1))$ & Right at root, Left at depth~1 \\
RR & $\eml(1,\eml(1,\eml(\cdot,\cdot)))$ & Right at both levels \\
LL & $\eml(\eml(\eml(\cdot,\cdot),1),1)$ & Left at both levels \\
\bottomrule
\end{tabular}
\end{center}

The fifth shape is \emph{balanced}: both root branches are full
subtrees, giving $\eml(\eml(\cdot,\cdot), \eml(\cdot,\cdot))$.
The LL chain produces double exponentiation ($e^{e^a}$), which
overflows on the training domain and is excluded from experiments.

Three architectures are tested, summarized in
Table~\ref{tab:archs}.  All three can represent any depth-3 EML tree
with $\{1, x, y\}$ at the leaves; one is strictly more expressive.
\textbf{Eq.\,6} is the architecture described in Equation~6 of
Odrzywolek~\citep{Odrzywolek2026eml}: at each internal node a
softmax selects among $\{1, x, f_{\text{child}}\}$, giving
variable~$x$ access to every level of the tree.  Because its input
set at internal nodes is a strict superset of V16's, Eq.\,6 is
strictly more expressive.
\textbf{V16} is variant~16 in the released source
code\footnote{Code: \url{https://github.com/aodrzywolek/EML}}
accompanying~\citet{Odrzywolek2026eml}: a sigmoid at each internal
node selects between $\{1, f_{\text{child}}\}$, and variables $x$
and $y$ enter only at the leaves through a 3-way softmax, making
them symmetric.
\textbf{Hybrid} uses V16's sigmoid at all internal nodes except the
root, where a 4-way softmax adds $x$ and $y$ as additional options;
the name reflects its combination of V16's body with partial variable
access at the root.

\begin{table}[t]
\centering
\caption{Architectures at depth~3.}
\label{tab:archs}
\begin{tabular}{@{}llrl@{}}
\toprule
Name & Internal selection & Params & Variable routing \\
\midrule
Eq.\,6  & softmax $\{1,x,f\}$ & 42 & $x$ at all nodes; $y$ leaves only \\
V16     & sigmoid $\{1,f\}$    & 38 & $x,y$ leaves only (symmetric) \\
Hybrid  & V16 + root 4-way     & 44 & $x,y$ root and leaves \\
\bottomrule
\end{tabular}
\end{table}


%% ===================================================================
\section{Methodology}
\label{sec:setup}
%% ===================================================================

To test whether the architecture--target interaction is specific to
EML or general to operators with gradient asymmetry, two additional
operators are introduced:
\begin{align}
  \sml(a,b) &= \sinh(a) - \arctan(b), \label{eq:sml} \\
  \rml(a,b) &= \arctan(a) - \sinh(b). \label{eq:rml}
\end{align}
SML shares EML's gradient direction: the left input receives
unbounded gradients ($\partial\sml/\partial a = \cosh a$, which
grows exponentially) while the right is bounded
($\partial\sml/\partial b = -1/(1+b^2)$).  RML reverses this
asymmetry.  Neither SML nor RML is known to be universal; target
trees are evaluated numerically.

Each node's inputs are weighted combinations of the available terms,
with weights controlled by learnable parameters.  Training minimizes
mean squared error using Adam~\citep{Kingma2015adam} (learning
rate 0.01) on a $21 \times 21$ grid over $[-3,3]^2$ (441~points;
targets with numerical overflow are filtered, leaving
$n_{\text{train}} \geq 378$).  A two-phase schedule is used:
a \emph{search phase} (6{,}000 iterations at temperature
$\tau = 2.5$, keeping softmax/sigmoid weights soft) followed by a
\emph{hardening phase} (2{,}000 iterations; $\tau$ anneals from
2.5 to 0.01 with entropy and binarity penalties, then snaps to
one-hot selections).  After hardening each node selects exactly one
input, producing a concrete expression tree.

Each combination of architecture, operator, tree shape, and leaf
assignment is tested with 256 independent trials
(64 seeds $\times$ 4 initializations) for V16 and Hybrid, or 64
trials (one initialization) for Eq.\,6; Eq.\,6 uses a single
initialization because the original paper specifies one, and adding
V16-style strategies to Eq.\,6 would confound the comparison.
Non-extreme rates (e.g., 17\%, 39.5\%) carry 95\%
Clopper--Pearson intervals of roughly $\pm 5$\,pp at $n = 256$.
A trial counts as
\emph{exact recovery} if the hardened tree matches the target
formula, verified on 500 random points (RMSE $< 10^{-6}$).
Replacing V16's initialization with Eq.\,6's
(randn $\times$ 0.1, no bias) yielded 0/64 recoveries, confirming
that the architecture differences reported below are not explained
by initialization.


%% ===================================================================
\section{Results}
\label{sec:results}
%% ===================================================================

\subsection{Architecture--Target Interaction}
\label{sec:results-interaction}

Table~\ref{tab:d3matrix} shows that
recovery depends on the combination of architecture and target, not
on either factor alone.  Eq.\,6 acts as a specialist: 100\% recovery
(64/64) on exactly two targets and 0\% (0/64) on three others.  V16
is a generalist with nonzero recovery on all four chain shapes
(16--99.6\%) and $x \leftrightarrow y$ symmetry from its leaf-only
variable design.  Hybrid is an RL specialist (95--96\% on RL targets,
near 0\% elsewhere).  No single architecture dominates: Eq.\,6
scores 100\% where V16 scores 17\%, yet V16 scores 99.6\% where
Eq.\,6 scores 0\%.  Expressiveness does not explain this pattern
because Eq.\,6 is strictly more expressive than V16 yet fails on
targets V16 recovers.  The difference lies in how each architecture's
variable routing interacts with the optimization landscape for a given
target.

\begin{table}[t]
\centering
\caption{EML recovery rates (\%, exact match).}
\label{tab:d3matrix}
\begin{tabular}{@{}llrrr@{}}
\toprule
Target & Shape & Eq.\,6 & V16 & Hybrid \\
\midrule
Paper (yx) & LR & \textbf{100} (64/64) & 17 & 0.4 \\
T1 (xy)    & RL & 0 (0/64) & \textbf{99.6} & \textbf{96} \\
T4 (xy)    & RR & 0 (0/64) & 39.5 & 0 \\
T8 (xy)    & LR & 1.6 & 16   & 0.4 \\
T1\_yx     & RL & 0   & ---  & \textbf{95.3} \\
T4\_yx     & RR & \textbf{100} (64/64) & --- & 0 \\
\bottomrule
\end{tabular}
\end{table}

\subsection{Operator Swap Reverses Shape Preferences}
\label{sec:results-operators}

Switching from EML to SML (Table~\ref{tab:sml},
Figure~\ref{fig:heatmap}) inverts Eq.\,6's shape preferences: its best shape under EML is LR (100\%), but
under SML it shifts to RL and RR (100\%) while LR drops to 0\%
(Fisher exact $p < 10^{-15}$).  Both operators amplify the left
input, so the reversal traces to differences in right-side
attenuation ($-1/(1+b^2)$ for SML versus $-1/b$ for EML), which
saturate differently on the $[-3,3]^2$ domain.  Despite the shape
reversal, qualitative profiles persist: Eq.\,6 remains a specialist,
V16 a generalist (83--100\% on all chains), and Hybrid an RL
specialist.  These profiles likely reflect structural properties of
variable routing rather than artifacts of one operator.

\begin{figure*}[t]
\centering
\includegraphics[width=0.82\textwidth]{fig_odrzywolek2026_heatmap.pdf}
\Description{Heatmap showing recovery rates for three operators across architectures and tree shapes.}
\caption{Recovery by architecture, shape, and operator (--- = inferred by symmetry).}
\label{fig:heatmap}
\end{figure*}

\begin{table}[t]
\centering
\caption{SML recovery rates (\%, exact match).}
\label{tab:sml}
\begin{tabular}{@{}llrrr@{}}
\toprule
Target & Shape & Eq.\,6 & V16 & Hybrid \\
\midrule
S\_LR (yx)    & LR & 0            & 84           & 0.8 \\
S\_LR\_xy     & LR & 0            & 83           & 0.8 \\
S\_RL (xy)    & RL & \textbf{100} & \textbf{100} & \textbf{100} \\
S\_RL\_yx     & RL & 0            & \textbf{100} & \textbf{100} \\
S\_RR (xy)    & RR & \textbf{100} & 99.6         & 89.8 \\
S\_RR\_yx     & RR & \textbf{100} & 98.8         & 88.7 \\
\bottomrule
\end{tabular}
\end{table}

V16 dominates under SML, recovering all chain targets at 83\% or
above.  Replacing EML with the right-amplified RML collapses all
rates: Eq.\,6 scores 0\% on every target, V16's best nominal
rate is 2.0\% (5/256, not significantly different from 0\%;
Fisher exact $p = 0.06$), and Hybrid's RL preference vanishes.  RR and balanced shapes
were not tested under RML because the collapse on chain targets,
which are easier to recover than balanced shapes, made further
experiments uninformative.


\subsection{Balanced Trees Are Universally Hard}
\label{sec:results-balanced}

The balanced depth-3 shape is tested with four leaf-placement
variants per architecture--operator combination
(Table~\ref{tab:balanced}).  Every architecture scores 0\% on every
balanced variant (0/3{,}776 trials; Clopper--Pearson 95\% upper
bound 0.10\%).  The failure persists across 6 training
configurations varying learning rate (0.01, 0.05), iteration budget
(6{,}000, 12{,}000), and temperature (1.0, 2.5): 0/192.  The
contrast with chain shapes is stark: V16 recovers 83--100\% of SML
chains yet scores exactly 0\% on balanced variants of the same
operator.  A plausible explanation is that balanced branching splits
gradient signal roughly equally between the two root subtrees,
preventing the concentration along a single path that drives chain
recovery; gradient measurements for balanced shapes were not
collected, and verifying this hypothesis is future work.

\begin{table}[t]
\centering
\caption{Balanced recovery: 0\% across all combinations.}
\label{tab:balanced}
\begin{tabular}{@{}lrrrr@{}}
\toprule
Operator & Eq.\,6 & V16 & Hybrid & Total \\
\midrule
SML (4 variants)  & 0/256  & 0/1024 & 0/1024 & 0/2{,}304 \\
EML (4 variants)  & 0/256  & 0/1024 & ---    & 0/1{,}280 \\
HP sensitivity    & ---    & 0/192  & ---    & 0/192 \\
\midrule
\textbf{Total}    &        &        &        & \textbf{0/3{,}776} \\
\bottomrule
\end{tabular}
\end{table}


\subsection{Gradient Measurements}
\label{sec:results-gradient}

To investigate the mechanism behind the architecture--target
interaction, gradient norms at leaf parameters are measured at
iteration~1000 across 10 seeds (Table~\ref{tab:gradients}).
For Eq.\,6, the ratio tracks recovery: when $x/y < 1$
(variable~$y$ receives stronger gradients), recovery is 100\%; when
$x/y > 1$ ($y$ is gradient-starved), recovery is 0\%.  Swapping EML
for SML on the same shape (LR) flips the ratio from 0.67 to 1.78,
consistent with the observed preference reversal.  V16 shows narrower
ratios across targets (0.77--2.00), consistent with its symmetric
design preventing either variable from dominating the gradient
landscape.

\begin{table}[t]
\centering
\caption{Leaf gradient ratio $\|\nabla_x\|/\|\nabla_y\|$ at iteration~1000
(10-seed mean $\pm$ s.e.).}
\label{tab:gradients}
\begin{tabular}{@{}llcrc@{}}
\toprule
Arch.\ & Target & Op.\ & $x/y$ ratio & Rec.\ \\
\midrule
Eq.\,6 & LR (yx) & EML & $0.67 \pm 0.03$ & 100\% \\
Eq.\,6 & RL (xy) & EML & $1.58 \pm 0.04$ & 0\% \\
Eq.\,6 & LR (yx) & SML & $1.78 \pm 0.12$ & 0\% \\
Eq.\,6 & RL (xy) & SML & $1.19 \pm 0.11$ & 100\% \\
\midrule
V16    & LR (yx) & EML & $0.77 \pm 0.05$ & 17\% \\
V16    & RL (xy) & EML & $2.00 \pm 0.08$ & 99.6\% \\
\bottomrule
\end{tabular}
\end{table}

The gradient trajectory during training (Figure~\ref{fig:gradient})
reveals the mechanism more clearly.  On the 100\% EML target the
ratio undergoes a brief reversal, from 1.45 at iteration~100 to
0.14 at iteration~500, then settling at 0.67 by iteration~1000;
this temporary window in which $y$ dominates correlates with
convergence to the correct formula.  On the 0\% target the ratio
remains above 1.15 at every measured iteration, and $y$ never
receives sufficient gradient signal.  These measurements establish a
gradient-level correlate of the architecture--target interaction; the
question of causation is addressed in Section~\ref{sec:limitations}.

\begin{figure}[t]
\centering
\includegraphics[width=\columnwidth]{fig_odrzywolek2026_gradient.pdf}
\Description{Line plot showing gradient ratio trajectories diverging for recovered vs unrecovered targets.}
\caption{Gradient ratio $\|\nabla_x\|/\|\nabla_y\|$ during training
(Eq.\,6, 10-seed mean $\pm$ s.e.).}
\label{fig:gradient}
\end{figure}


%% ===================================================================
\section{Discussion}
\label{sec:discussion}
%% ===================================================================

Selecting a single architecture for all targets, as current
gradient-based SR methods do~\citep{Sahoo2018eql, Dong2024eeql,
Odrzywolek2026eml}, can yield 0\% recovery on targets that a
different architecture recovers at 100\%.  Neural architecture search
(NAS) also studies architecture--task interaction, but NAS
architectures typically differ in
expressiveness~\citep{Elsken2019nas}; here all architectures share
the same target language, so recovery differences arise purely from
optimization landscape geometry, a condition rarely met in benchmarks
where methods differ simultaneously in grammar, search, and
architecture~\citep{LaCava2021srbench}.

The cross-operator results point to four interacting factors that
determine recovery: variable routing, tree shape, leaf assignment,
and gradient direction.  Four predictions were formulated before the
RML and balanced experiments.  One was confirmed (V16 preferences
invert under RML); three were falsified (balanced trees were
predicted to vary by architecture but proved universally hard;
Eq.\,6 and Hybrid were predicted to reverse preferences under RML
but collapsed instead).  Gradient direction interacts with variable
routing more strongly than a simple reordering model predicts.

\subsection{Limitations}
\label{sec:limitations}

All three operators share the subtraction form $f(a) - g(b)$;
whether the interaction extends to multiplicative or compositional
operators is untested.  The full three-architecture comparison exists
only at depth~3; limited depth-4 experiments with V16 show the
generalist profile collapses (RL drops from 99.6\% to 0.8\%), but a
complete depth-4 matrix remains future work.

The controlled-expressiveness design requires all targets to lie
within the EML grammar, preventing direct comparison with standard SR
benchmarks (SRBench, Feynman equations) that would require
exponential-depth EML trees.  SML targets have narrower output range
than EML targets on $[-3,3]^2$, which may partly explain V16's
higher SML rates (83--100\% versus 16--99.6\%).  The shape-specific
patterns (Eq.\,6 under SML: LR 0\%, RL 100\%) cannot be explained by
conditioning alone, since better conditioning would raise all shapes
uniformly.

The gradient measurements establish correlation, not causation.  A
mechanistic account via loss-landscape analysis would strengthen the
framework but is beyond the scope of this study.


%% ===================================================================
\section{Conclusion}
\label{sec:conclusion}
%% ===================================================================

The optimization landscape, not expressiveness alone, determines what
gradient-based symbolic regression recovers.  This motivates two
practical directions: multi-architecture ensembles that run several
structures in parallel and select the best result, and
architecture-aware search that adapts variable routing to the target
during optimization.  The universal failure on balanced trees points
to a deeper constraint that warrants direct investigation via
loss-landscape analysis.  Extending the interaction analysis to
multiplicative operators, deeper trees, and standard SR benchmarks
would test how far these findings generalize.


%% ===================================================================
%% Bibliography (pre-compiled .bbl inlined for arXiv)
%%% -*-BibTeX-*-
%%% Do NOT edit. File created by BibTeX with style
%%% ACM-Reference-Format-Journals [18-Jan-2012].

\begin{thebibliography}{14}

%%% ====================================================================
%%% NOTE TO THE USER: you can override these defaults by providing
%%% customized versions of any of these macros before the \bibliography
%%% command.  Each of them MUST provide its own final punctuation,
%%% except for \shownote{} and \showURL{}.  The latter two
%%% do not use final punctuation, in order to avoid confusing it with
%%% the Web address.
%%%
%%% To suppress output of a particular field, define its macro to expand
%%% to an empty string, or better, \unskip, like this:
%%%
%%% \newcommand{\showURL}[1]{\unskip}   % LaTeX syntax
%%%
%%% \def \showURL #1{\unskip}           % plain TeX syntax
%%%
%%% ====================================================================

\ifx \showCODEN    \undefined \def \showCODEN     #1{\unskip}     \fi
\ifx \showISBNx    \undefined \def \showISBNx     #1{\unskip}     \fi
\ifx \showISBNxiii \undefined \def \showISBNxiii  #1{\unskip}     \fi
\ifx \showISSN     \undefined \def \showISSN      #1{\unskip}     \fi
\ifx \showLCCN     \undefined \def \showLCCN      #1{\unskip}     \fi
\ifx \shownote     \undefined \def \shownote      #1{#1}          \fi
\ifx \showarticletitle \undefined \def \showarticletitle #1{#1}   \fi
\ifx \showURL      \undefined \def \showURL       {\relax}        \fi
% The following commands are used for tagged output and should be
% invisible to TeX
\providecommand\bibfield[2]{#2}
\providecommand\bibinfo[2]{#2}
\providecommand\natexlab[1]{#1}
\providecommand\showeprint[2][]{arXiv:#2}

\bibitem[Biggio et~al\mbox{.}(2021)]%
        {Biggio2021nsrs}
\bibfield{author}{\bibinfo{person}{Luca Biggio}, \bibinfo{person}{Tommaso
  Bendinelli}, \bibinfo{person}{Alexander Neitz}, \bibinfo{person}{Aurelien
  Lucchi}, {and} \bibinfo{person}{Giambattista Parascandolo}.}
  \bibinfo{year}{2021}\natexlab{}.
\newblock \showarticletitle{Neural Symbolic Regression that Scales}. In
  \bibinfo{booktitle}{\emph{Proceedings of the 38th International Conference on
  Machine Learning (ICML)}}. \bibinfo{pages}{936--945}.
\newblock


\bibitem[Cava et~al\mbox{.}(2021)]%
        {LaCava2021srbench}
\bibfield{author}{\bibinfo{person}{William~La Cava}, \bibinfo{person}{Patryk
  Orzechowski}, \bibinfo{person}{Bogdan Burlacu},
  \bibinfo{person}{Fabr{\'\i}cio~Olivetti de Fran{\c{c}}a},
  \bibinfo{person}{Marco Virgolin}, \bibinfo{person}{Ying Jin},
  \bibinfo{person}{Michael Kommenda}, {and} \bibinfo{person}{Jason~H. Moore}.}
  \bibinfo{year}{2021}\natexlab{}.
\newblock \showarticletitle{Contemporary Symbolic Regression Methods and their
  Relative Performance}. In \bibinfo{booktitle}{\emph{Advances in Neural
  Information Processing Systems (NeurIPS)}}, Vol.~\bibinfo{volume}{34}.
\newblock


\bibitem[Dong and Zhong(2024)]%
        {Dong2024eeql}
\bibfield{author}{\bibinfo{person}{Hao Dong} {and} \bibinfo{person}{Jiangjun
  Zhong}.} \bibinfo{year}{2024}\natexlab{}.
\newblock \showarticletitle{Extended Equation Learner Solves Symbolic
  Regression Tasks with Evolutionary Components}.
\newblock \bibinfo{journal}{\emph{IEEE Transactions on Evolutionary
  Computation}} (\bibinfo{year}{2024}).
\newblock


\bibitem[Dong and Zhong(2025)]%
        {Dong2025survey}
\bibfield{author}{\bibinfo{person}{Hao Dong} {and} \bibinfo{person}{Jiangjun
  Zhong}.} \bibinfo{year}{2025}\natexlab{}.
\newblock \showarticletitle{A Survey on Symbolic Regression}.
\newblock \bibinfo{journal}{\emph{arXiv preprint arXiv:2211.10873}}
  (\bibinfo{year}{2025}).
\newblock


\bibitem[Elsken et~al\mbox{.}(2019)]%
        {Elsken2019nas}
\bibfield{author}{\bibinfo{person}{Thomas Elsken}, \bibinfo{person}{Jan~Hendrik
  Metzen}, {and} \bibinfo{person}{Frank Hutter}.}
  \bibinfo{year}{2019}\natexlab{}.
\newblock \showarticletitle{Neural Architecture Search: A Survey}.
\newblock \bibinfo{journal}{\emph{Journal of Machine Learning Research}}
  \bibinfo{volume}{20}, \bibinfo{number}{55} (\bibinfo{year}{2019}),
  \bibinfo{pages}{1--21}.
\newblock


\bibitem[Ipek(2026a)]%
        {Ipek2026battery}
\bibfield{author}{\bibinfo{person}{Sinan Ipek}.}
  \bibinfo{year}{2026}\natexlab{a}.
\newblock \showarticletitle{Evaluating the Exp-Minus-Log Sheffer Operator for
  Battery Characterization}.
\newblock \bibinfo{journal}{\emph{arXiv preprint arXiv:2604.13873}}
  (\bibinfo{year}{2026}).
\newblock


\bibitem[Ipek(2026b)]%
        {Ipek2026hardware}
\bibfield{author}{\bibinfo{person}{Sinan Ipek}.}
  \bibinfo{year}{2026}\natexlab{b}.
\newblock \showarticletitle{Hardware-Efficient Neuro-Symbolic Networks with the
  Exp-Minus-Log Operator}.
\newblock \bibinfo{journal}{\emph{arXiv preprint arXiv:2604.13871}}
  (\bibinfo{year}{2026}).
\newblock


\bibitem[Kingma and Ba(2015)]%
        {Kingma2015adam}
\bibfield{author}{\bibinfo{person}{Diederik~P. Kingma} {and}
  \bibinfo{person}{Jimmy Ba}.} \bibinfo{year}{2015}\natexlab{}.
\newblock \showarticletitle{Adam: A Method for Stochastic Optimization}.
\newblock \bibinfo{journal}{\emph{Proceedings of the International Conference
  on Learning Representations (ICLR)}} (\bibinfo{year}{2015}).
\newblock


\bibitem[Landajuela et~al\mbox{.}(2022)]%
        {Landajuela2022udsr}
\bibfield{author}{\bibinfo{person}{Mikel Landajuela},
  \bibinfo{person}{Chak~Shing Lee}, \bibinfo{person}{Jiachen Yang},
  \bibinfo{person}{Ruben Glatt}, \bibinfo{person}{Claudio~P. Santiago},
  \bibinfo{person}{Ignacio Aravena}, \bibinfo{person}{Terrell Mundhenk},
  \bibinfo{person}{Garrett Mulcahy}, {and} \bibinfo{person}{Brenden~K.
  Petersen}.} \bibinfo{year}{2022}\natexlab{}.
\newblock \showarticletitle{A Unified Framework for Deep Symbolic Regression}.
  In \bibinfo{booktitle}{\emph{Advances in Neural Information Processing
  Systems (NeurIPS)}}, Vol.~\bibinfo{volume}{35}.
\newblock


\bibitem[Li et~al\mbox{.}(2025)]%
        {Li2025metasymnet}
\bibfield{author}{\bibinfo{person}{Wenqiang Li}, \bibinfo{person}{Weijun Li},
  \bibinfo{person}{Lina Yu}, \bibinfo{person}{Min Wu}, \bibinfo{person}{Linjun
  Sun}, \bibinfo{person}{Jingyi Liu}, \bibinfo{person}{Yanjie Li}, {and}
  \bibinfo{person}{Songsong Tian}.} \bibinfo{year}{2025}\natexlab{}.
\newblock \showarticletitle{{MetaSymNet}: A Dynamic Symbolic Regression Network
  Capable of Evolving into Arbitrary Formulas}. In
  \bibinfo{booktitle}{\emph{Proceedings of the AAAI Conference on Artificial
  Intelligence}}.
\newblock


\bibitem[Makke and Chawla(2024)]%
        {Makke2024survey}
\bibfield{author}{\bibinfo{person}{Nour Makke} {and} \bibinfo{person}{Sanjay
  Chawla}.} \bibinfo{year}{2024}\natexlab{}.
\newblock \showarticletitle{Interpretable Scientific Discovery with Symbolic
  Regression: A Review}.
\newblock \bibinfo{journal}{\emph{Artificial Intelligence Review}}
  \bibinfo{volume}{57}, \bibinfo{number}{2} (\bibinfo{year}{2024}).
\newblock


\bibitem[Odrzywolek(2026)]%
        {Odrzywolek2026eml}
\bibfield{author}{\bibinfo{person}{Andrzej Odrzywolek}.}
  \bibinfo{year}{2026}\natexlab{}.
\newblock \showarticletitle{All Elementary Functions from a Single Operator}.
\newblock \bibinfo{journal}{\emph{arXiv preprint arXiv:2603.21852v2}}
  (\bibinfo{year}{2026}).
\newblock
\newblock
\shownote{Institute of Theoretical Physics, Jagiellonian University}.


\bibitem[Petersen et~al\mbox{.}(2021)]%
        {Petersen2021dsr}
\bibfield{author}{\bibinfo{person}{Brenden~K. Petersen}, \bibinfo{person}{Mikel
  Landajuela}, \bibinfo{person}{T.~Nathan Mundhenk},
  \bibinfo{person}{Claudio~P. Santiago}, \bibinfo{person}{Soo~K. Kim}, {and}
  \bibinfo{person}{Joanne~T. Kim}.} \bibinfo{year}{2021}\natexlab{}.
\newblock \showarticletitle{Deep Symbolic Regression: Recovering Mathematical
  Expressions from Data via Risk-Seeking Policy Gradients}. In
  \bibinfo{booktitle}{\emph{Proceedings of the International Conference on
  Learning Representations (ICLR)}}.
\newblock


\bibitem[Sahoo et~al\mbox{.}(2018)]%
        {Sahoo2018eql}
\bibfield{author}{\bibinfo{person}{Subham Sahoo}, \bibinfo{person}{Christoph
  Lampert}, {and} \bibinfo{person}{Georg Martius}.}
  \bibinfo{year}{2018}\natexlab{}.
\newblock \showarticletitle{Learning Equations for Extrapolation and Control}.
  In \bibinfo{booktitle}{\emph{Proceedings of the 35th International Conference
  on Machine Learning (ICML)}}. \bibinfo{pages}{4442--4450}.
\newblock


\end{thebibliography}

\end{document}
